\section{A finite split for bounded posterior densities}
\label{l:bounded-split:section}

We partition a bounded-density posterior into finitely many components,
each supported in a dyadic cell and satisfying a bound on mass at every
smaller scale.  Revealing the component has an entropy cost.  An exact
change of reference from the parent cell to the selected cell compares
this cost with the information supplied by the cell's depth.  We first
construct the partition, then prove the block estimate, and finally combine them
into a finite-block information bound and a learner lower bound.

Put \(n=d-1\), and use the closed cube
\[
 Q=[-1/(2\sqrt n),1/(2\sqrt n)]^n,\qquad
 s(z)=(z,\sqrt{1-\|z\|^2}).
\]
Let \(\tau\) be its uniform probability and assign dyadic boundaries
consistently.  The geometric proof of Lemma~\ref{ld:cube-geometry}
applies on this closure: \(s\) is \(2\)-Lipschitz, a depth-\(h\)
cell has side \(2^{-h}/\sqrt n\) and mass \(2^{-nh}\), and the
first \(n\) graph coordinates recover \(z\).  For a cell \(A\), let
\(\tau_A\) be uniform probability on it.
Throughout this section set
\[
 k=\lfloor n/64\rfloor,\qquad J=2k,\qquad p=\lfloor n/8\rfloor,
\]
and take \(n\) sufficiently large that \(k\ge1\) and \(p\ge2\).
Here \(k\) is the block length, \(J\) is the number of rows in two
stacked blocks, and \(p\) is the projection moment order.
Constants denoted by \(C,c\) are positive and absolute and may change
between estimates.  The independent-draw, span-distance projection
calculation is related to that of Sharan, Sidford, and Valiant
\cite{SSV2019}; the exact-label density and block estimates are proved
below.

\subsection{A finite partition and its entropy identity}

\begin{lemma}[Bounded-density split]\label{l:bounded-split:decomposition}
Let \(\mu\) be a probability supported on a dyadic cell \(A\),
with bounded density relative to \(\tau_A\).  There is a finite
measurable partition into positive-probability labels \(i\), with
probabilities \(p_i\), descendant cells \(A_i\) of relative depths
\(j_i\), and numbers \(q_i\in(0,1]\).  Writing
\(\mu_i=\mu(\,\cdot\mid i)\) and \(Q_i=\log(1/q_i)\), each
\(\mu_i\) is supported on \(A_i\) and satisfies
\begin{equation}\label{l:bounded-split:spread}
 \mu_i(D)\le q_i^{-1}2^{-n\ell/2}
\end{equation}
for every relative-depth-\(\ell\) descendant \(D\) of \(A_i\), and
\begin{align}
 -\log p_i&\le\tfrac12n(\log2)j_i+Q_i,
 &\mathbb E_iQ_i&\le1+\mathbb E_i j_i.
 \label{l:bounded-split:label-cost}
\end{align}
\begin{equation}\label{l:bounded-split:entropy-identity}
 \mathbb E_iD(\mu_i\Vert\tau_{A_i})
 =D(\mu\Vert\tau_A)+H(i)-n(\log2)\mathbb E_i j_i.
\end{equation}
The descendant bound includes \(\ell=0\), and \(\mathbb E_i\)
means expectation with weights \(p_i\).
\end{lemma}

\begin{proof}
At a node cell \(C\), use the original law \(\mu\) conditioned on
\(C\).  Select all inclusion-maximal strict dyadic descendants \(D\)
whose conditional mass exceeds \(2^{-n\ell/2}\), where \(\ell\)
is depth relative to \(C\).  They are disjoint: two dyadic cells are
disjoint or nested, and nesting would contradict maximality.
A point in a selected cell moves to that cell.  Every other point stops
in the remainder of the current cell.  Repeat at the selected cells.
The label records the terminal node and its remainder.

Let \(K=\|d\mu/d\tau_A\|_\infty<\infty\); necessarily \(K\ge1\).
Along a nontrivial path reaching relative depth \(j\), the product of
the selected conditional masses exceeds \(2^{-nj/2}\).
No complement decision occurs before stopping, so reaching the path
cell is exactly membership in it.  The same product is its original
\(\mu\)-mass, which is at most \(K2^{-nj}\).  Hence
\[
 j<\frac{2\log K}{n\log2}.
\]
This uniformly bounds every nontrivial path depth and leaves finitely
many grid cells.  At an individual node, its bounded conditional
density likewise rules out heavy descendants at unbounded depth,
so each selected inclusion-maximal descendant exists.  The
construction is therefore a finite tree and defines a finite partition.

Let \(q_i\) be the conditional probability of the terminal remainder
at its node, omitting zero-probability remainders.  A strict descendant
that exceeded its threshold before the remainder decision is contained
in a selected cell and has zero remainder mass.  Every other
descendant has its conditional mass increased by at most \(q_i^{-1}\).
This proves \eqref{l:bounded-split:spread}; the case \(\ell=0\)
is immediate.  The product of masses on the path, multiplied by
\(q_i\), is \(p_i\), so its lower bound proves
\(-\log p_i\le\tfrac12n(\log2)j_i+Q_i\).

For the expected remainder cost, a node reached with probability \(r\)
and having remainder probability \(q\) contributes
\(rq\log(1/q)\) to \(\mathbb E_iQ_i\).  This is at most \(r\),
since \(q\log(1/q)\le1\), with value zero at \(q=0\).
Summing over the finite tree bounds \(\mathbb E_iQ_i\) by the
expected number of visited nodes.  Each move increases depth by at
least one, so this number is at most \(1+\mathbb E_i j_i\).

On the part with label \(i\), the density relative to \(\tau_{A_i}\)
is the original density relative to \(\tau_A\), multiplied by
\(p_i^{-1}2^{-nj_i}\).  Taking its logarithm, integrating on that
part, and averaging over \(i\) proves
\eqref{l:bounded-split:entropy-identity}.  All terms are finite
because the partition is finite and each conditional density is bounded.
\end{proof}

The entropy identity credits \(n(\log2)\mathbb E_i j_i\) for the
change of cell.  Naming the component uses at most half of this credit,
plus \(\mathbb E_iQ_i\).  The remaining task is to bound the next
block's information in component \(i\) by \(Cn+CQ_i\); the bound
\(\mathbb E_iQ_i\le1+\mathbb E_i j_i\) will then pay for its
dependence on the remainder probability.

\subsection{A centered exact projection density}

Squaring a state likelihood introduces two independent \(k\)-row
matrices at the same signal.  We will therefore need a density estimate
for their \(J=2k\) stacked rows.  The \(p\) independent signal draws
in the proof below serve a different purpose: they expand the \(p\)th
moment of that density.

Consider any probability \(\mu\) on a depth-\(h\) cell \(A\) that
satisfies \eqref{l:bounded-split:spread} with \(q\in(0,1]\).
Put \(Q_*= \log(1/q)\), \(R=2^{-h}\), and let \(s_0\) be the
graph image of the cell center.  The cell has radius \(R/2\), so
\(\|s(z)-s_0\|\le R\).  The following estimates use only this
dyadic mass hypothesis; the bounded-density premise was used to
construct the finite split.

\begin{lemma}[Local density for exact labels]
\label{l:bounded-split:density}
Let \(\mu\) be a probability on a depth-\(h\) dyadic cell \(A\)
such that, for some \(q\in(0,1]\), every relative-depth-\(\ell\)
descendant \(D\) satisfies \(\mu(D)\le q^{-1}2^{-n\ell/2}\).
Put \(Q_*=\log(1/q)\), \(R=2^{-h}\), and
\(s_0=s(\operatorname{center}A)\), with \(J,p\) as fixed above.
Let \(X'\) have \(J\) independent standard Gaussian rows, independently
of \(z\sim\mu\), and put \(b=X'(s(z)-s_0)\).
On the open region
\(\Omega=\{b\in\mathbb R^J:\|b\|<3\sqrt J R\}\), the joint law
of \((X',b)\) has a density \(f\) relative to \(dP_{X'}\,db\).
If \(h_R\) is the density of \(N(0,R^2I_J)\), then
\begin{equation}\label{l:bounded-split:reference-norm}
 \left\|\mathbf1_\Omega f/h_R\right\|_
 {L^p(dP_{X'}h_R(b)\,db)}
 \le C^J e^{C(J/n)Q_*},
\end{equation}
where \(\mathbf1_\Omega f/h_R\) is extended by zero outside
\(\Omega\).
\end{lemma}

\begin{proof}
We first use the dyadic mass bound to control inverse Gram determinants
of independent centered signal vectors.  These determinants bound the
moment of a smoothed projection density.  We then remove the smoothing
and compare the exact density with the Gaussian reference.

For every \(y\in\mathbb R^d\) and \(0<r\le1\),
\begin{equation}\label{l:bounded-split:ball}
 \Pr_\mu\{\|s(z)-y\|\le rR\}\le C^ne^{Q_*}r^{n/2}.
\end{equation}
Indeed, choose \(\ell\ge0\) with
\(2^{-\ell-1}<r\le2^{-\ell}\).  A relative-depth-\(\ell\)
cell meeting the projection of this ball onto the first \(n\)
coordinates lies in the concentric ball of radius
\(2\,2^{-\ell}R\).  Its side is \(2^{-\ell}R/\sqrt n\).
The unit-ball volume estimate in Lemma~\ref{ld:cube-geometry}
therefore bounds the number of such disjoint cells by \(C^n\).
Multiply by their mass bound \(e^{Q_*}2^{-n\ell/2}\) and use
\(2^{-\ell}<2r\).

If \(U\subseteq\mathbb R^d\) is a linear subspace of dimension
\(m\le n/8\), a maximal \(rR\)-separated set in its radius-\(R\)
ball has at most \((1+2/r)^m\) points.  Its radius-\(rR\)
balls cover that ball, by maximality.  Because \(\|s-s_0\|\le R\),
ambient balls of radius \(2rR\) about the translated centers cover
the event \(\operatorname{dist}(s-s_0,U)\le rR\).
Using \eqref{l:bounded-split:ball}, and the trivial mass bound for
\(r>1/2\), gives
\begin{equation}\label{l:bounded-split:tube}
 \Pr_\mu\{\operatorname{dist}(s-s_0,U)\le rR\}
 \le C^ne^{Q_*}r^{n/2-m}\qquad(0<r\le1).
\end{equation}

For \(m\le p-1\), put \(a=n/2-m\ge3n/8\) and
\(D_U=R^{-1}\operatorname{dist}(s-s_0,U)\).
Then \(0<D_U\le1\) almost surely.  Let
\(\mathcal A=C^ne^{Q_*}\ge1\).  Since \(J\le a/2\), layer-cake
integration of \(\Pr(D_U<r)\le\min\{1,\mathcal A r^a\}\), split at
\(\mathcal A^{-1/a}\), gives
\begin{equation}\label{l:bounded-split:inverse}
 \begin{split}
 \mathbb E_\mu D_U^{-J}
 &=1+J\int_0^1r^{-J-1}\Pr(D_U<r)\,dr\\
 &\le\frac a{a-J}\mathcal A^{J/a}
 \le2\mathcal A^{J/a}\le(Ce^{CQ_*/n})^J.
 \end{split}
\end{equation}
For independent \(z_1,\ldots,z_p\sim\mu\), put \(s_i=s(z_i)\)
and let \(G\) be the Gram matrix of \(e_i=s_i-s_0\).  Conditional on its
predecessors, each \(e_i\) has positive distance from their span
almost surely by \eqref{l:bounded-split:tube}.  Gram--Schmidt and
successive conditioning in \eqref{l:bounded-split:inverse} yield
\begin{equation}\label{l:bounded-split:gram}
 \mathbb E(\det G)^{-J/2}
 \le(Ce^{CQ_*/n}/R)^{Jp}.
\end{equation}

Add an independent uniform vector in \([-t,t]^J\), \(t>0\), to
the exact label \(b\).  Its joint density is
\[
 f_t(X',b)=(2t)^{-J}\mathbb E_\mu
   \mathbf1_{\{\|X'(s-s_0)-b\|_\infty\le t\}}.
\]
Expand its integer \(p\)th power using independent
\(e_1,\ldots,e_p\).  For fixed \(b\), Tonelli's theorem and
independence of the \(J\) rows give
\[
 \int f_t(X',b)^p\,dP_{X'}
 =(2t)^{-Jp}\mathbb E_{e_1,\ldots,e_p}
 \prod_{j=1}^J
 \Pr_x\{|\langle x,e_i\rangle-b_j|\le t\text{ for all }i\le p\}.
\]
For almost every tuple the inner-product vector for one row is
Gaussian with nonsingular covariance \(G\).  Its density is bounded
by \((2\pi)^{-p/2}(\det G)^{-1/2}\).  Each box probability is
at most \((2t)^p\) times this supremum.  The \(2t\) powers cancel,
and \eqref{l:bounded-split:gram} gives, uniformly in \(t,b\),
\[
 \int f_t(X',b)^p\,dP_{X'}
 \le(Ce^{CQ_*/n}/R)^{Jp}.
\]
Integrating on \(\Omega\) therefore bounds the local \(p\)th
power integral by
\(|\Omega|(Ce^{CQ_*/n}/R)^{Jp}\).

For a continuous compactly supported test function on the product
of the matrix space and \(\Omega\), its integral against the
jittered law converges to its integral against the exact law as
\(t\downarrow0\).  Extend the test by zero outside a compact
subset of \(\Omega\) and use bounded convergence.
H\"older bounds these limiting integrals by the local \(L^p\)
bound times the test's \(L^{p/(p-1)}(dP_{X'}\,db)\) norm.
Density of such tests and \(L^p\) duality give a density for the
exact law on this open region, with
\begin{equation}\label{l:bounded-split:local-lp}
 \int_\Omega\int f(X',b)^p\,dP_{X'}\,db
 \le|\Omega|(Ce^{CQ_*/n}/R)^{Jp}.
\end{equation}
Equality on the tests identifies the represented measure with the
restricted exact law.  This asserts a density on \(\Omega\), which
is the only region used below.

On \(\Omega\), the Gaussian reference satisfies
\(h_R(b)\ge(CR)^{-J}\).  Also the ball-volume estimate gives
\(|\Omega|\le(CR)^J\).  In the integral of \(f^ph_R^{1-p}\),
the powers of \(R\) from \eqref{l:bounded-split:local-lp} are
\[
 R^{-Jp}R^JR^{J(p-1)}=1.
\]
Taking the \(p\)th root gives
\eqref{l:bounded-split:reference-norm}.  The auxiliary jitter has
been removed before the exact observation rule is applied.
\end{proof}

\subsection{Two copies of the block and the state information}

\begin{lemma}[Diffuse block information]\label{l:bounded-split:block}
Let \(\mu\) be a probability on a depth-\(h\) dyadic cell \(A\)
such that, for some \(q\in(0,1]\), every relative-depth-\(\ell\) descendant
\(D\) satisfies \(\mu(D)\le q^{-1}2^{-n\ell/2}\), and put
\(Q_*=\log(1/q)\).  Suppose \(\log W\le n^2\), with
\(k=\lfloor n/64\rfloor\) as above.
From a fixed entering state, let \(u\)
be any variable with at most \(W\) values obtained by a measurable
probability kernel from a \(k\)-row standard Gaussian matrix \(X\)
and exact labels \(Xs(z)\), where \(X\) is independent of \(z\sim\mu\).
Then
\begin{equation}\label{l:bounded-split:block-information}
 I_\mu(z;u)\le Cn+CQ_*.
\end{equation}
\end{lemma}

\begin{proof}
Put \(R=2^{-h}\) and \(s_0=s(\operatorname{center}A)\).
Write the kernel probabilities as \(g_u(X,Y)\in[0,1]\), with
\(\sum_ug_u(X,Y)=1\) for every matrix-label pair \((X,Y)\).  Define
\[
 \mathcal G=\{\|X(s-s_0)\|\le2\sqrt k R\},\qquad
 K_u(s)=\mathbb E_X[\mathbf1_{\mathcal G}g_u(X,Xs)],
\]
and define a comparison probability vector by
\[
 \rho_u=\mathbb E_{X,Z}g_u(X,Xs_0+RZ),
 \qquad Z\sim N(0,I_k)\text{ independently}.
\]
In \(\mathbb E_\mu K_u(s)^2\), use two independent block matrices
at the same signal \(s\).  Their stack has \(J=2k\) rows.
When both blocks satisfy \(\mathcal G\), the stacked residual norm
is at most \(2\sqrt J R<3\sqrt J R\), inside \(\Omega\).
Apply Lemma~\ref{l:bounded-split:density} and H\"older's inequality
to this exact joint law.  Under the Gaussian reference, the two
matrices and the two residual vectors are independent between blocks.
The product of the two \(g_u\)'s lies in \([0,1]\) and has
reference expectation \(\rho_u^2\).  Dropping the good indicators
in its conjugate norm gives
\begin{equation}\label{l:bounded-split:square}
 \mathbb E_\mu K_u(s)^2
 \le C^J e^{C(J/n)Q_*}\rho_u^{2(1-1/p)}.
\end{equation}
Thus \(\rho_u=0\) implies \(K_u=0\) almost surely.

For every fixed \(s\), the probability of \(\mathcal G^c\) is
at most
\[
 e^{-k}2^{k/2}=e^{-ck}=:\delta<1
\]
for an absolute \(c>0\), by the Gaussian exponential moment at
\(1/4\).  Put \(g_0(s)=\sum_uK_u(s)\ge1-\delta\) and
\(P_g=\mathbb E_\mu g_0(s)\).  Including the good indicator
with \(u\) can only increase information.  The indicator costs at
most \(\log2\), and the weighted bad conditional information costs
at most \(\delta\log W\).

The weighted good conditional information is at most
\begin{equation}\label{l:bounded-split:good-kl}
 \mathbb E_\mu\sum_uK_u(s)
       \log\frac{K_u(s)}{g_0(s)\rho_u}.
\end{equation}
Indeed, replacing the actual conditional marginal
\(\Pr(u\mid\mathcal G)\) by \(\rho\) adds
\(P_gD(\Pr(u\mid\mathcal G)\Vert\rho)\ge0\).
Zero states cause no issue by \eqref{l:bounded-split:square}.
The contribution of \(-\log g_0(s)\) is at most
\(\log(1/(1-\delta))\).  For the remaining term, Jensen's inequality
under the probability measure with weights
\(\mu(ds)K_u(s)/P_g\) gives
\[
 \begin{split}
 \mathbb E_\mu\sum_uK_u(s)\log\frac{K_u(s)}{\rho_u}
 &\le P_g\log\left(\frac1{P_g}
      \sum_{u:\rho_u>0}\frac{\mathbb E_\mu K_u(s)^2}{\rho_u}\right)\\
 &\le P_g\log\left(
      \frac{C^J e^{C(J/n)Q_*}W^{2/p}}{P_g}\right).
 \end{split}
\]
The last inequality uses \eqref{l:bounded-split:square} and
\(\sum_u\rho_u^{1-2/p}\le W^{2/p}\), by concavity for \(p\ge2\).
All sums are finite.  Combining the terms gives
\[
 I_\mu(z;u)\le Cn+CQ_*+\frac2p\log W+e^{-ck}\log W.
\]
Since \(k,p\) are proportional to \(n\) and \(\log W\le n^2\),
the last two terms are \(O(n)\).  This proves the lemma.
\end{proof}

\subsection{The finite-block information bound}

The split identity and the block estimate can now be combined without
an accuracy or stopping assumption.  We keep every split label in the
analysis history, so its entropy must be included in the bound.

\begin{proposition}[Information after finite splitting]
\label{l:bounded-split:iteration}
Let \(z\sim\tau\), and let \(B\ge0\) and \(W\ge1\) be integers
with \(\log W\le n^2\).  Independently of \(z\), draw independent
standard Gaussian \(k\times d\) matrices
\(X_1,\ldots,X_B\), with \(k\) as fixed above.  Consider a
computation with a fixed initial state \(u_0\) and at most \(W\) states
at each block boundary.  At block \(b\),
the outgoing state \(u_b\) is an everywhere-defined Borel function of
\((u_{b-1},X_b,X_bs(z))\).

Before each block one can reveal a finite partition label \(i_b\),
determined by \(z\) and the preceding analysis history, such that the
complete history
\[
 \mathcal H=(i_1,u_1,\ldots,i_B,u_B)
\]
has finite range and satisfies
\(I(z;\mathcal H)\le4C'nB\) for an absolute constant \(C'>0\)
and all sufficiently large \(n\).  The labels indicate nested dyadic
cells containing \(z\); they do not alter the computation's state rules.
\end{proposition}

\begin{proof}
Before each block, apply Lemma~\ref{l:bounded-split:decomposition}
to the law conditional on the analysis transcript, and record its
label.  Also record the outgoing boundary state after the block.
Initially the cell is the whole cube and the law is \(\tau\).
Every positive-probability posterior has bounded density relative to
its current cell's uniform law: splitting preserves this property,
and conditioning on a state multiplies the density by a likelihood
in \([0,1]\) and divides by a positive state probability.
Every split is finite, so each stage has finitely many histories.
Future Gaussian rows remain independent of the signal and these
histories, which depend only on the signal and past blocks.
The block from a known entering state is a measurable kernel of
its exact data.  Lemma~\ref{l:bounded-split:block} therefore applies
in each selected component.

For a current cell of depth \(h\) and conditional law \(\mu\),
define
\begin{equation}\label{l:bounded-split:potential}
 \Phi=\tfrac14n(\log2)h+D(\mu\Vert\tau_A).
\end{equation}
It is nonnegative.  By \eqref{l:bounded-split:entropy-identity},
the expected change under a split is
\[
 H(i)-\tfrac34n(\log2)\mathbb E_i j_i
 \le-\tfrac14n(\log2)\mathbb E_i j_i+\mathbb E_iQ_i.
\]
Within component \(i\), conditioning on its outgoing state keeps
the cell fixed and increases the expected relative entropy by exactly
\(I_{\mu_i}(z;u)\).  This follows by multiplying the density by
\(\Pr(u\mid z,i)/\Pr(u\mid i)\), taking logarithms, and averaging.
The block lemma bounds this increase by \(Cn+CQ_i\).
Thus the combined conditional drift is at most
\begin{equation}\label{l:bounded-split:drift}
 \begin{split}
 Cn-\tfrac14n(\log2)\mathbb E_i j_i+(C+1)\mathbb E_iQ_i
 &\le Cn+(C+1)
   +(C+1-\tfrac14n\log2)\mathbb E_i j_i\\
 &\le C'n
 \end{split}
\end{equation}
for large \(n\), using \(\mathbb E_iQ_i\le1+\mathbb E_i j_i\).
Only this expected bound on \(Q_i\) is used.

The initial potential is zero, so after \(B\) blocks
\(\mathbb E\Phi\le C'nB\).  For the final transcript
\(\mathcal H\), its posterior in a depth-\(h\) cell satisfies
\[
 D(P(z\mid\mathcal H)\Vert\tau)
 =n(\log2)h+D(P(z\mid\mathcal H)\Vert\tau_A)
 \le4\Phi.
\]
Averaging over the finite history gives
\begin{equation}\label{l:bounded-split:total-information}
 I(z;\mathcal H)\le4C'nB.
\end{equation}
\end{proof}

\subsection{Application to the learner}

Use the finite-state model of Section~\ref{l:model}, with
\(M(d)=o(d^2)\), integer horizon \(T\), and every-signal angular
success at least \(2/3\) at \(0<\epsilon\le1/10\).
We will prove \(T\ge c d\log(1/\epsilon)\), with an absolute
constant and a dimension threshold that may depend on the memory
sequence.  First record the route's separate sample bound for the
original randomized learner:
\begin{equation}\label{l:bounded-split:linear}
 T\ge d/16.
\end{equation}
Take a uniform full-sphere signal and pregenerate all \(T\) rows,
including unused rows.  Let \(V\) be their span.  Conditional on
the rows and \(P_VS\), the residual has mean zero; this remains true
after conditioning on the learner's independent random tape.
All labels and the output are functions of the conditioned variables.
Consequently
\[
 \mathbb E\langle S,\widehat S\rangle
 \le\mathbb E\|P_VS\|
 \le\sqrt{\mathbb E\|P_VS\|^2}
 \le\sqrt{T/d}.
\]
The last inequality uses isotropy and independence of the rows and
signal.  Success gives
\[
 \mathbb E\langle S,\widehat S\rangle
 \ge(2/3)\cos(1/10)-1/3>1/4,
\]
which proves \eqref{l:bounded-split:linear}.

Put \(L=\log(1/\epsilon)\).  In the cube experiment, first fix
shared rule randomness with success greater than \(1/2\).
The graph derivative has full rank, so Lemmas~\ref{l:prior-ac}
and~\ref{l:borel} give Borel representatives for that fixed choice.
Then fix the remaining transition and output randomness, together with
the initial state, by averaging.  Conditional on the shared seed these
choices are independent of the signal and rows, by the joint
independence in Section~\ref{l:model}.  This gives a deterministic
program of cube success at least \(1/2\).  Its terminal state and
stopping index determine its output, giving at most \((T+1)2^M\)
possible outputs.  Lemma~\ref{ld:cube-geometry} bounds the cube success
mass of each fixed output by \((5\epsilon)^n\).  Hence
\begin{equation}\label{l:bounded-split:output-count}
 n(L-\log5)\le M\log2+\log(2(T+1)).
\end{equation}

We may restrict to \(T<dL\), since the complementary case already
has the required lower bound.  For \(L\ge\log10\), put
\(a_0=\log2/\log10>0\); then \(L-\log5\ge a_0L\).
Also \(\log(2(T+1))\le C+\log d+\log L\le C+\log d+L\).
Thus
\[
 (a_0n-1)L\le M\log2+C+\log d,
\]
so \(L=o(d)\).  Pad the deterministic program after an early stop,
retaining its terminal state and stopping index and ignoring later
samples.  The width
\begin{equation}\label{l:bounded-split:width}
 W=(T+2)2^M
\end{equation}
is enough for this padded program.  Since \(T<dL\), \(L=o(d)\),
and \(M=o(d^2)\), it satisfies \(\log W=o(n^2)\), hence
\(\log W\le n^2\) eventually.  Use \(B=\lceil T/k\rceil\)
blocks, adding ignored fresh rows to complete the last block.

The fixed initial state, independent blocks, and Borel transition rules
now satisfy Proposition~\ref{l:bounded-split:iteration}.  Let
\(\mathcal H\) be its augmented history.  The padded state retains
the terminal state and stopping index, so the output is a function of
\(\mathcal H\).

Under the joint law, the output's first \(n\) coordinates are within
\(\epsilon\) of \(z\)
with probability at least \(1/2\); under the product law the
probability is at most \((5\epsilon)^n\) by the coordinate-ball
volume calculation in Lemma~\ref{ld:cube-geometry}.
Equation~\eqref{l:binary} yields
\begin{equation}\label{l:bounded-split:accuracy}
 I(z;\mathcal H)\ge\tfrac12n(L-\log5)-\log2\ge cnL
\end{equation}
for sufficiently large \(n\), uniformly for \(L\ge\log10\).
Combining \eqref{l:bounded-split:total-information} and
\eqref{l:bounded-split:accuracy} gives \(B\ge c_1L\).
Since \(B\le T/k+1\),
\[
 T\ge k(c_1L-1).
\]
For \(L\ge2/c_1\), this is \(\Omega(dL)\).
For \(\log10\le L<2/c_1\), the separate bound
\eqref{l:bounded-split:linear} gives the same order.
The previously excluded case \(T\ge dL\) is immediate.  This proves
the worst-signal precision lower bound for every stated accuracy
sequence, with absolute constants and a dimension threshold that may
depend on the memory sequence.
